\scanpage{12}
épinglage, où $\Pi_0$ est le ``\unsure{$n$-polyg.|polyg.|$n$-poly.} combinatoire type''
% slanted
\margin{(i.e. \inserted{``le''} \deleted{\unsure{tout|deux}} \deleted{\unsure{le|l'un} \illegible{1} épinglé} pol. comb. épinglé -- i.e. muni d'un \uline{repère}\ldots)}
\uline{Cor. 2} Soient $P=(E,S,\omega)$, $P'=(E',S',\omega')$ deux
polygônes réguliers isomorphes, \deleted{Alors} le
premier étant épinglé \inserted{par un repère $r$}. Alors l'application
\[
\unsure{\Theta|\theta}\unsure{\longmapsto|\longrightarrow}\Theta(r)\,:\quad \operatorname{Isom}(P,P')\longrightarrow\operatorname{Rep}(E',S',\omega')
\]
est bijective.

\uline{Scholie} On peut dire encore\,: dans tte
classe d'isomorphie de $n$-polygônes
réguliers, les $n$-polygônes réguliers \uline{épinglés}
sont déterminés à \uline{isom} \uuline{unique} près.
[Idem bien sûr ds le cas purement
combinatoire]. \inserted{Bien sûr, les classes d'iso des uns et des autres se correspondent 1-1\ldots}

\uline{NB} Il y aura donc la question ``d'écrire''\
``\uline{le}'' espace affine $E$ dans lequel se
place $P$, en termes du type d'isomorphie
(i.e. de $n$ et d'un
``module'' $\alpha\in k$ convenable), \deleted{et l'angle
au sommet} et ceci fait, ``d'écrire''\
également \uline{les} \deleted{\unsure{côtés}} \inserted{sommets} $s_i$ ($0\leq i\unsure{\leq|<} n-1$).
% archive sheet number: 11

